\chapter{Tinjauan Pustaka dan Dasar Teori}

\section{Tinjauan Pustaka}
\subsection{Penelitian Komputasional pada Jaringan Sanad Hadis}
\subsubsection{Analisis Jaringan Sanad}

Penelitian komputasional pada sanad hadis berkembang dari sekadar pembangunan dataset menuju pemodelan jaringan berbasis graf dan penerapan \textit{Social Network Analysis} (SNA) untuk mengukur posisi tiap perawi.

Mghari \textit{et al.} (2022) membangun \textit{Sanadset 650K}, korpus hadis yang memuat 650.986 entri dari 926 kitab berbahasa Arab menggunakan \textit{web scraping} dan \textit{web crawling} berbasis Python (\textit{Beautiful Soup} dan \textit{Selenium}). Setiap entri menyertakan teks hadis lengkap, matan, daftar perawi, dan urutan rantai sanad. Dataset ini dipublikasikan di repositori Mendeley Data untuk mendukung klasifikasi hadis (sahih/lemah), pengenalan perawi, pemodelan linguistik Arabic NLP, dan pembangunan graf perawi \cite{MGHARI2022108540}. Skala dan cakupannya yang mencapai ratusan kitab menjadikan Sanadset 650K sumber data paling representatif untuk penelitian jaringan transmisi sanad.

\begin{sloppypar}
Farooqi \textit{et al.} (2024) memperkenalkan \textit{Multi-IsnadSet} (MIS), dataset yang merepresentasikan jaringan transmisi Sahih Muslim sebagai graf berarah dengan 2.092 \textit{node} (perawi) dan 77.797 \textit{edge} \cite{FAROOQI2024110439}. Untuk membangun dan menganalisis jaringan tersebut, peneliti menggunakan NetworkX, Neo4j, Gephi, dan Cytoscape. Studi ini menunjukkan bahwa jaringan sanad dapat dimodelkan sebagai \textit{knowledge graph} di Neo4j dan dianalisis dengan metrik SNA (\textit{degree}, \textit{betweenness}, \textit{eigenvector centrality}) untuk mengidentifikasi perawi berpengaruh \cite{FAROOQI2024110439}. Namun, visualisasi yang dihasilkan bersifat statis dan bergantung pada perangkat lunak desktop, sehingga pengguna umum tidak dapat mengaksesnya tanpa instalasi khusus.
\end{sloppypar}

Shuaib \textit{et al.} (2022) menerapkan \textit{trophic analysis}, metode yang diadaptasi dari ekologi makanan, terhadap \textit{Hadith Social Network} dari corpus Gawāmi' al-Kalim yang memuat 447.205 hadis dan 49.309 perawi unik dengan 294.303 edge berarah, mencakup sembilan abad sejarah transmisi \cite{Shuaib2022}. Jaringan sanad diperlakukan sebagai sistem difusi informasi hierarkis, dengan Nabi Muhammad saw. sebagai \textit{node} sumber (\textit{basal node}) dan tiap lapisan transmisi mencerminkan jarak temporal dari sumber. Hasilnya menunjukkan bahwa jaringan sanad bersifat hierarkis dan sekuensial, dan analisis berbasis graf mengungkap informasi temporal yang tidak tercatat dalam teks historis \cite{Shuaib2022}. Shuaib \textit{et al.} menyebutkan potensi penggabungan analisis trofik dengan metrik sentralitas SNA untuk mengidentifikasi perawi berpengaruh pada periode tertentu, sebuah arah yang belum diwujudkan dalam antarmuka yang dapat dieksplorasi pengguna.

\subsubsection{Knowledge Graph untuk Representasi Jaringan Perawi}

Mahmoud \textit{et al.} (2024) memperkenalkan \textit{NarratorsKG} sebagai \textit{knowledge graph} pertama yang didedikasikan untuk perawi hadis \cite{Mahmoud2024}. KG ini dibangun sebagai graf berarah G = (V, E) dengan dua jenis \textit{node}: \textit{node} perawi (\textit{narrator nodes}) dan \textit{node} bentuk penampilan nama (\textit{appearance form nodes}). NarratorsKG memuat 79.256 \textit{node}, terdiri dari 18.298 \textit{node} perawi dan 60.958 \textit{node} bentuk penampilan, serta 197.183 \textit{edge} yang mencakup 70.415 \textit{edge} perawi--penampilan dan 126.768 \textit{edge} hubungan "meriwayatkan dari" \cite{Mahmoud2024}.

KG ini digunakan bersama pendekatan \textit{graph query embedding} dua tahap: tahap pertama menghasilkan \textit{embedding} dari struktur graf sanad, lalu tahap kedua menggunakan \textit{AraBERT} sebagai \textit{reranker} untuk mengidentifikasi perawi berdasarkan informasi biografis (nama lengkap, teknonym, tahun wafat, negara, dan silsilah). Diuji pada AR-Sanad 280K-v2, pendekatan ini mencapai akurasi 95,2\% pada \textit{validation set} \cite{Mahmoud2024}. KG berstruktur ini meningkatkan akurasi identifikasi perawi dibanding pendekatan berbasis teks biasa, tetapi NarratorsKG dirancang sebagai sumber daya internal untuk tugas NLP, bukan sistem yang dapat dieksplorasi pengguna akhir.

\subsubsection{Social Network Analysis dan Metrik Sentralitas}

\begin{sloppypar}
Studi SNA pada jaringan perawi hadis menunjukkan pola yang berulang pada berbagai korpus. Pada Sahih Bukhari dan Sahih Muslim, \textit{degree centrality} mengidentifikasi perawi paling aktif seperti Abu Hurayra dan Anas bin Malik, \textit{betweenness centrality} mengungkap posisi "jembatan" tokoh seperti Imam Malik ibn Anas, dan \textit{PageRank} menempatkan Az-Zuhri serta Sufyan bin Uyaynah sebagai perawi paling berpengaruh secara struktural \cite{Mghari2024Narrator2Vec}. Tiap metrik menangkap dimensi pengaruh yang berbeda, sehingga penggunaannya bersama memberikan gambaran posisi perawi yang lebih lengkap.

Mghari \textit{et al.} (2024) menawarkan alternatif melalui \textit{Narrator2Vec}, adaptasi \textit{Skip-Gram Word2Vec} yang melatih representasi vektor 100 dimensi dari nama perawi pada lebih dari 650.000 hadis. Model ini mengidentifikasi nama-nama berbeda dari perawi yang sama dan memprediksi perawi yang hilang dalam rantai terputus \cite{Mghari2024Narrator2Vec}. Narrator2Vec tidak menghasilkan skor kuantitatif per perawi yang dapat ditampilkan dan dibandingkan pengguna, sehingga metrik sentralitas lebih sesuai untuk eksplorasi interaktif jaringan perawi.
\end{sloppypar}

\subsubsection{Otomatisasi Analisis Sanad dengan Pendekatan Machine Learning}

Hendawi \textit{et al.} (2024) mengusulkan metode berbasis \textit{machine learning} untuk memverifikasi kesinambungan sanad pada hadis terputus (\textit{disconnected hadith}), dengan memanfaatkan data biografis perawi, meliputi tanggal lahir dan wafat, pergerakan geografis, serta penilaian ulama dari ilmu Jarh wa Ta'dil sebagai fitur prediksi \cite{10974901}. \textit{Random Forest Classifier} dan \textit{HistGradientBoosting Classifier} mencapai akurasi tertinggi, masing-masing 96,55\% dan 93,16\%, melampaui \textit{AraBERT} (\textit{underfitting} karena dataset kecil) dan \textit{ChatGPT} yang hanya mencapai 77,23\% \cite{10974901}. Hasil ini menunjukkan bahwa tugas verifikasi sanad lebih bergantung pada data biografis spesifik daripada kemampuan pemahaman bahasa umum, dan memperkuat kebutuhan representasi data perawi yang terstruktur sebagai landasan pembangunan knowledge graph.

\begin{table}[H]
\centering
\caption{Perbandingan Studi Dataset Hadis}
\label{tab:ringkasan-pustaka}
\renewcommand{\arraystretch}{1.3}
\resizebox{\textwidth}{!}{%
\begin{tabular}{|>{\raggedright\arraybackslash}p{3cm}|>{\raggedright\arraybackslash}p{2.5cm}|>{\raggedright\arraybackslash}p{3cm}|>{\raggedright\arraybackslash}p{3.5cm}|>{\raggedright\arraybackslash}p{2.5cm}|>{\raggedright\arraybackslash}p{4cm}|>{\raggedright\arraybackslash}p{4cm}|}
\hline
\multicolumn{1}{|c|}{\textbf{Peneliti (Tahun)}} & \multicolumn{1}{c|}{\textbf{Dataset Utama}} & \multicolumn{1}{c|}{\textbf{Tujuan}} & \multicolumn{1}{c|}{\textbf{Metode}} & \multicolumn{1}{c|}{\textbf{Alat/Platform}} & \multicolumn{1}{c|}{\textbf{Hasil Utama}} & \multicolumn{1}{c|}{\textbf{Keterbatasan}} \\
\hline
Mghari, Bouras \& El Hibaoui (2022) \cite{MGHARI2022108540} &
Sanadset 650K (926 kitab) &
Membangun dataset sanad skala besar terbuka &
\textit{Web scraping} \& \textit{crawling} + \textit{tagging} NLP (SANAD/MATN/NAR) &
Python (\textit{Beautiful Soup}, \textit{Selenium}) &
Korpus 650.986 entri dengan matan, sanad, dan perawi terpisah; fondasi penelitian sanad komputasional &
Hanya dataset mentah; tanpa visualisasi graf maupun analisis jaringan \\
\hline
Mghari, Bouras \& El Hibaoui (2024) \cite{Mghari2024Narrator2Vec} &
Sanadset 650K ($\pm$650 rb hadis) &
Membangun representasi vektor perawi &
Word embedding Skip-Gram Word2Vec (Narrator2Vec, 100 dimensi) &
Python (Gensim), t-SNE &
Mampu mengidentifikasi nama berbeda dari perawi yang sama \& memprediksi perawi hilang; akurasi top-1 tinggi &
Tidak menghasilkan skor sentralitas kuantitatif per perawi yang dapat dibandingkan pengguna \\
\hline
Farooqi, Malick, Shaikh \& Akhunzada (2024) \cite{FAROOQI2024110439} &
Multi-IsnadSet (Sahih Muslim, 2.092 perawi) &
Menyediakan dataset graf multi-isnad \& knowledge graph &
Graf berarah + knowledge graph; menyarankan SNA &
Neo4j, Gephi, Cytoscape, NetworkX &
2.092 node \& 77.797 edge; KG dibangun di Neo4j; sediakan basis untuk SNA &
Visualisasi statis; bergantung perangkat desktop (Gephi/Cytoscape); tidak interaktif/berbasis web \\
\hline
Shuaib, Syed, Halawi \& Saquib (2022) \cite{Shuaib2022} &
Gaw\={a}mi' al-Kalim (447.205 hadis, 49.309 perawi) &
Merekonstruksi informasi temporal dari jaringan &
\textit{Trophic analysis} (adaptasi ekologi) &
Python, NetworkX &
\textit{Trophic level} berkorelasi dengan tahun wafat perawi; jaringan terbukti hierarkis \& sekuensial &
Tidak memakai metrik sentralitas; tanpa antarmuka eksplorasi interaktif \\
\hline
Mahmoud, Nabil, Saif \& Torki (2024) \cite{Mahmoud2024} &
AR-Sanad 280K-v2 (279.625 sanad) &
Disambiguasi perawi dalam sanad &
Knowledge graph (NarratorsKG) + graph query embedding + AraBERT &
NarratorsKG, AraBERT, Python &
KG perawi pertama (79.256 node, 197.183 edge); akurasi disambiguasi 95,2\% &
KG untuk NLP internal; tidak tersedia sebagai antarmuka publik yang dapat dieksplorasi \\
\hline
Hendawi, Torki, Elmasry \& Khalafallah (2024) \cite{10974901} &
Data biografis perawi (Mousoua al-Hadith) &
Verifikasi kesinambungan sanad terputus &
Machine learning (TF-IDF + RF, HistGradBoost) &
Python (Scikit-learn), AraBERT &
Random Forest 96,55\% \& HistGradBoost 93,16\%; ungguli ChatGPT (77,23\%) &
Dataset kecil (503 paragraf); tidak mencakup semua alasan terputusnya sanad \\
\hline
\textbf{Penelitian Ini (2026)} &
Sanadset 650K (926 kitab) &
Membangun \textit{Knowledge Graph} jaringan keilmuan Islam serta mengidentifikasi peran dan pengaruh perawi secara kuantitatif &
Pra-pemrosesan data, pembangunan KG sebagai graf berarah berbobot, SNA dengan 5 metrik sentralitas &
Python (NetworkX), React, Supabase &
KG dengan 177.047 \textit{node} \& 776.798 \textit{edge}; \textit{dashboard} web interaktif yang dapat diakses publik tanpa instalasi &
\textit{Closeness} \& \textit{betweenness centrality} dihitung pada subgraf terkuat; data bersifat statis; belum menilai kesahihan sanad \\
\hline
\end{tabular}%
}
\end{table}

\subsection{Visualisasi Perawi Dalam Bentuk Graf}

Beberapa penelitian merepresentasikan rantai perawi (isnad/sanad) sebagai graf berarah, dengan tiap perawi sebagai simpul (\textit{node}) dan tiap relasi periwayatan sebagai sisi (\textit{directed edge}). Model ini memungkinkan isnād yang dibaca linear dievaluasi sebagai jaringan berlapis yang mencerminkan hubungan historis antar perawi \cite{ahmad2025hadith}. Saraçoğlu menerapkan analisis jaringan sosial terhadap isnād dalam kitab Kit\={a}bu't-\d{T}abaqa\={a}ti'l-Kab\={i}r karya Ibn Sa'd, dengan mengonversi 2.258 rantai isnād dari hampir 2.000 perawi menjadi data terstruktur berbasis pasangan \textit{source--target} \cite{saraccougluvisualization}. Visualisasi dilakukan menggunakan perangkat lunak seperti \textit{Pajek}, \textit{Gephi}, dan \textit{Isnalyser}, di mana perawi berperan sebagai aktor dan relasi periwayatan melalui lafaz seperti akhbaran\={a}, \d{h}addathan\={a}, dan 'an berperan sebagai ikatan. Penelitian tersebut juga menyoroti tantangan visualisasi jaringan historis jangka panjang, khususnya dalam menata kronologi aliran informasi dan menangani panjangnya nama-nama perawi dalam bahasa Arab \cite{saraccougluvisualization}.

Pendekatan serupa dikembangkan melalui sistem berbasis web bernama HadithNet, yang mengotomatisasi identifikasi nama perawi dan menghasilkan rantai perawi kronologis pada teks hadis terjemahan Melayu menggunakan Django, MySQL, dan pustaka Graphviz \cite{alias2024hadithnet}. Sistem ini memanfaatkan teknik \textit{preprocessing}, \textit{POS tagging}, dan \textit{Named Entity Recognition} (NER), serta mencapai akurasi tinggi dalam identifikasi nama perawi dengan presisi 99,37\%, meskipun masih menghadapi kendala dalam memvisualisasikan rantai perawi yang kompleks dan bercabang banyak. Penelitian ini juga merangkum beberapa pendekatan visualisasi terdahulu, seperti graf jaringan jarang (\textit{sparse network graph}) untuk menghubungkan sepuluh perawi paling berpengaruh, serta jaringan berskala besar yang dibangun menggunakan pustaka Python seperti NetworkX, yang menunjukkan bahwa kompleksitas keterhubungan menjadi tantangan utama dalam interpretasi visual \cite{alias2024hadithnet}.

Ahmad \textit{et al.}\ \cite{ahmad2025hadith} menunjukkan bahwa transformasi isnād menjadi graf berarah memungkinkan penerapan metrik analisis jaringan, seperti \textit{connectivity}, \textit{chain depth}, deteksi \textit{cluster}, \textit{bridge narrators}, dan sentralitas (\textit{betweenness}, \textit{eigenvector}, \textit{closeness}), untuk mengidentifikasi pola, klaster regional, dan anomali struktural dalam rantai periwayatan. Representasi ini memungkinkan sarjana melihat keseluruhan jaringan isnād sekaligus, membandingkan jalur kanonik dan non-kanonik, serta menemukan titik lemah struktural yang tidak terdeteksi melalui pembacaan manual \cite{ahmad2025hadith}. Sementara itu, sistem KASHAF memperkenalkan visualisasi alir (\textit{flow-visualization}) berbasis diagram Sankey untuk menganalisis lebih dari 60.000 hadis dari koleksi Sembilan Kitab, dengan menampilkan grade perawi, hubungan jalur periwayatan, serta inkonsistensi dalam rantai sanad secara interaktif \cite{shafie2021kashaf}.

Tabel~\ref{tab:visualisasi-perawi} merangkum perbandingan keempat studi di atas berdasarkan sistem yang digunakan, metode dan teknologi, dataset, hasil evaluasi, serta keterbatasannya.

\begin{table}[H]
\centering
\caption{Perbandingan Studi Visualisasi Perawi Hadis}
\label{tab:visualisasi-perawi}
\renewcommand{\arraystretch}{1.3}
\resizebox{\textwidth}{!}{%
\begin{tabular}{|>{\raggedright\arraybackslash}p{3cm}|>{\raggedright\arraybackslash}p{3cm}|>{\raggedright\arraybackslash}p{4cm}|>{\raggedright\arraybackslash}p{4cm}|>{\raggedright\arraybackslash}p{4cm}|>{\raggedright\arraybackslash}p{3.5cm}|}
\hline
\multicolumn{1}{|c|}{\textbf{Studi}} & \multicolumn{1}{c|}{\textbf{Sistem/Alat}} & \multicolumn{1}{c|}{\textbf{Metode \& Teknologi}} & \multicolumn{1}{c|}{\textbf{Dataset/Uji Coba}} & \multicolumn{1}{c|}{\textbf{Hasil Evaluasi}} & \multicolumn{1}{c|}{\textbf{Keterbatasan}} \\
\hline
Sara\c{c}o\u{g}lu (2024) \cite{saraccougluvisualization} &
Pajek, Gephi, Isnalyser &
Analisis jaringan sosial; konversi isnād ke pasangan \textit{source--target}; jaringan berarah &
2.258 rantai isnād, $\pm$2.000 perawi dari Kit\={a}bu't-\d{T}abaqa\={a}ti'l-Kab\={i}r (Ibn Sa'd) &
Visualisasi jaringan kronologis (contoh: 162 isnād, 30 node, 28 edge) &
Sulit menata kronologi \& panjang nama Arab pada jaringan jangka panjang \\
\hline
Alias \textit{et al.}\ (2024) \cite{alias2024hadithnet} &
Sistem web (Django, MySQL, Graphviz) &
NER berbasis aturan, POS tagging, \textit{preprocessing}; visualisasi graf Graphviz &
1.000 teks hadis Sahih Bukhari terjemahan Melayu &
Presisi 99,37\% (nama); F1 rantai tunggal 96,21\%, rantai ganda 53,33\% &
Lemah pada rantai perawi kompleks/bercabang banyak \\
\hline
Ahmad \textit{et al.}\ (2025) \cite{ahmad2025hadith} &
Model ML (GNN, \textit{transformer}, klasifikasi hibrida) &
Representasi graf \& sekuens; metrik sentralitas, \textit{clustering}, \textit{bridge}; disambiguasi perawi &
AR-Sanad 280K, Sanadset 650K, Multi-IsnadSet (Sahih Muslim) + verifikasi rijal &
Akurasi disambiguasi $>$85\% (terdokumentasi), $<$60\% (langka); selaras penilaian klasik &
Tak menilai aspek moral; bergantung kekayaan data; perlu pengawasan ahli \\
\hline
KASHAF -- Shafie (2021) \cite{shafie2021kashaf} &
Mesin pencari hadis berbasis \textit{knowledge graph} &
\textit{Querying knowledge graph}; klasifikasi \textit{takhreej} dengan AraBERT; visualisasi alir Sankey &
60.000+ hadis koleksi Sembilan Kitab; 200K pasangan untuk klasifikasi &
\textit{Recall} \& presisi 0,9 pada klasifikasi \textit{takhreej}; visualisasi Sankey interaktif &
Fokus pada \textit{retrieval} \& visualisasi, bukan penilaian \textit{grade} hadis otomatis \\
\hline
\textbf{Penelitian Ini (2026)} &
\textit{Dashboard} web (React, Supabase, HTML5 Canvas) &
\textit{Knowledge graph} + SNA (5 metrik sentralitas); visualisasi graf \textit{force-directed} interaktif &
Sanadset 650K (487.451 rantai sanad bersih) &
177.047 \textit{node}, 776.798 \textit{edge}; pengujian \textit{black-box} 14 skenario, tingkat keberhasilan 100\% &
\textit{Closeness}/\textit{betweenness centrality} pada subgraf terkuat; belum menilai kesahihan sanad; data belum diperbarui otomatis \\
\hline
\end{tabular}%
}
\end{table}

\section{Dasar Teori}


\subsection{Hadis}
Secara etimologi, kata hadis bermakna al-jadīd (yang baru), al-qarīb (yang dekat), dan al-khabar (berita), sementara menurut istilah ahli hadis, hadis adalah segala sesuatu yang disandarkan kepada Nabi Muhammad saw., baik berupa perkataan, perbuatan, maupun ketetapan (taqrir), sehingga terbagi menjadi tiga jenis, yaitu hadis qauli, fi'li, dan taqrir \cite{daulay2023hadis}. 
Kedudukan hadis sangat penting dalam pendidikan karena Nabi sendiri menaruh perhatian besar terhadap persoalan pendidikan, dan generasi yang berkualitas hanya dapat terbentuk melalui pemahaman agama yang bersumber pada Al-Qur'an dan hadis \cite{daulay2023hadis}.

Sebagai sumber hukum Islam, hadis menempati posisi kedua setelah Al-Qur'an, sehingga ketaatan terhadap keduanya tidak dapat dipisahkan dan seorang muslim wajib merujuk pada keduanya secara bersamaan dalam memahami syariat \cite{siregar2024alquran}. 
Dalam hal ini hadis berperan sebagai penjelas (bayan) terhadap Al-Qur'an yang bersifat umum dan global, dengan fungsi yang mencakup penguatan (taqrir), perincian (tafsir), penetapan hukum baru (tasyri'), serta penghapusan ketentuan terdahulu (nasakh) menurut sebagian ulama \cite{siregar2024alquran}.

Mengingat tidak semua hadis dapat dijadikan dasar hukum, para ulama mengklasifikasikan hadis berdasarkan kualitasnya menjadi tiga kategori, yaitu hadis sahih, hasan, dan dhaif \cite{pratiwi2024klasifikasi}. 
Hadis sahih adalah hadis yang sanadnya bersambung, diriwayatkan oleh perawi yang adil dan dhabith, serta bebas dari syadz dan illat, sehingga dapat dijadikan hujjah dalam seluruh aspek hukum Islam. 
Hadis hasan memiliki kualitas yang menyerupai hadis sahih, namun tingkat ketelitian perawinya sedikit lebih rendah, sehingga tetap dapat dijadikan landasan hukum kecuali pada persoalan aqidah. 
Adapun hadis dhaif merupakan hadis yang mengandung kelemahan pada sanad, perawi, atau matannya, sehingga penggunaannya terbatas pada keutamaan amal (fadhilah amal) dengan syarat tertentu dan tidak dapat dijadikan dasar penetapan hukum syar'i \cite{pratiwi2024klasifikasi}.

\subsection{Sanad}
Secara etimologis, kata sanad berasal dari bahasa Arab yang bermakna sandaran atau pegangan, sedangkan secara terminologis sanad didefinisikan sebagai rangkaian para perawi yang menyampaikan dan menukilkan matan hadis dari sumber pertamanya hingga sampai kepada periwayat terakhir \cite{Mustang_Abubakar_2025}. 
Bersama matan dan rawi, sanad merupakan salah satu dari tiga unsur utama penyusun hadis, dan menempati posisi yang sangat sentral karena berfungsi sebagai jalur transmisi yang menghubungkan suatu hadis dengan sumber aslinya, yaitu Rasulullah saw.
Oleh karena itu, keberadaan sanad menjadi sangat krusial dalam menentukan diterima atau ditolaknya suatu hadis, sebagaimana ungkapan Abdullah ibn al-Mubarak yang menegaskan bahwa sanad merupakan bagian dari agama, dan tanpanya setiap orang dapat menyampaikan pendapatnya sesuka hati \cite{Mustang_Abubakar_2025}.

Sebuah hadis hanya dapat diterima apabila sanadnya bersambung (ittiṣāl al-sanad) dan diriwayatkan oleh para perawi yang bersifat tsiqah, yang mencakup dua unsur, yaitu 'adil (integritas moral keagamaan) dan ḍābiṭ (ketelitian serta kekuatan hafalan), serta terhindar dari syadz dan illat. 
Kriteria keshahihan sanad ini telah dirumuskan oleh para ulama dan secara konsisten digunakan lintas generasi sebagai fondasi metodologis dalam kritik sanad (naqd al-sanad), yakni upaya menilai keotentikan hadis melalui penelusuran kualitas para perawi yang membentuk rangkaian periwayatannya \cite{Mustang_Abubakar_2025}.

\subsection{Perawi}
Perawi (rawi) adalah individu yang terlibat dalam proses penerimaan dan penyampaian hadis dari satu generasi ke generasi berikutnya, dan kualitasnya menjadi salah satu faktor utama yang menentukan diterima atau ditolaknya suatu hadis \cite{Syafii2022Ketsiqohan}. 
Tinggi rendahnya sifat adil dan dhabit para perawi menyebabkan kuat atau lemahnya kedudukan suatu hadis, sehingga perawi yang memenuhi kedua sifat tersebut disebut tsiqah (terpercaya), sedangkan lawannya disebut dha'if, yaitu perawi yang cacat dalam hal keadilan dan kedhabitan. 
Sifat adil mencerminkan integritas keagamaan dan kepribadian baik (muru'ah) seorang perawi, sementara dhabit menunjukkan ketelitian serta kekuatan hafalan dan dokumentasinya; keduanya harus terpenuhi sekaligus agar seorang perawi memperoleh predikat tsiqah. Untuk menilai kredibilitas ini, para ulama mengembangkan cabang ilmu jarh wa al-ta'dil sebagai instrumen untuk mengetahui mana hadis yang sahih dan mana yang dhaif \cite{Syafii2022Ketsiqohan}.

Dalam praktiknya, periwayatan hadis merupakan proses penerimaan (tahammul) dan penyampaian (ada') hadis dengan menyebutkan sanad secara jelas, yang dijaga keakuratannya melalui hafalan, penulisan, dan pengamalan. Periwayatan ini terbagi menjadi dua bentuk, yaitu periwayatan bil lafadz, yang mempertahankan redaksi asli hadis sebagaimana diterima dari Rasulullah saw. dan dipandang lebih kuat, serta periwayatan bil makna, yang membolehkan penggunaan redaksi berbeda selama substansi maknanya tetap terjaga \cite{hatta2026teknik}. Agar periwayatannya dapat diterima, seorang perawi harus memenuhi sejumlah syarat, yakni beragama Islam, mukallaf (baligh dan berakal), bersifat adil, memiliki kedhabitan, serta meriwayatkan melalui sanad yang bersambung, tidak syadz, dan terbebas dari cacat tersembunyi (illah qadihah) \cite{hatta2026teknik}.

Adapun metode penerimaan dan penyampaian hadis terdiri atas delapan cara utama, yaitu sama' (mendengar langsung dari guru), qira'ah 'ala al-syaikh (membacakan di hadapan guru), ijazah, munawalah, mukatabah, i'lam, wasiat, dan wijadah. 
Di antara metode tersebut, sama' dipandang sebagai yang paling kuat karena terjadi transmisi langsung antara guru dan murid, dan keberagaman metode ini menunjukkan bahwa tradisi periwayatan hadis berkembang secara sistematis dalam menjaga keotentikan ajaran Islam \cite{hatta2026teknik}.

\subsection{Graf}
Graf $G$ didefinisikan sebagai pasangan dari dua himpunan berhingga, yaitu himpunan titik dan himpunan sisi, sebagaimana dirumuskan pada Persamaan~\eqref{eq:graf}.
\begin{equation}
G = (V, E)
\label{eq:graf}
\end{equation}
dengan $V$ menyatakan himpunan titik dan $E$ menyatakan himpunan sisi.
Titik (\textit{vertex} atau \textit{node}) merupakan komponen dasar penyusun graf, sedangkan sisi (\textit{edge}) merupakan elemen yang menghubungkan dua buah titik.
Suatu graf memungkinkan untuk tidak memiliki sisi sama sekali, namun setidaknya harus memuat satu titik. Himpunan titik dinyatakan dengan $V(G) = \{v_1, v_2, v_3, \dots, v_n\}$, sedangkan himpunan sisi dinyatakan dengan $E(G) = \{e_1, e_2, e_3, \dots, e_n\}$. 
Banyaknya titik pada graf $G$ disebut sebagai \textit{order} dan dilambangkan dengan $|V(G)|$, sementara banyaknya sisi disebut \textit{size} dan dilambangkan dengan $|E(G)|$. 
Dua titik $u$ dan $v$ pada graf $G$ dikatakan saling bertetangga (\textit{adjacent}) apabila keduanya dihubungkan oleh suatu sisi $e$, yaitu $e = uv$, dan banyaknya sisi yang bersisian dengan suatu titik $u$ disebut sebagai derajat (\textit{degree}) titik tersebut yang dinotasikan dengan $d(u)$ \cite{Lestari_Kristiana_Prihandini_Alfarisi_Setiawan_Adawiyah_2024}.

\subsection{\textit{Knowledge Graph}}
\textit{Knowledge graph} merupakan representasi data berbentuk grafik yang dirancang untuk menghimpun dan menyajikan informasi secara terstruktur. Dalam struktur ini, setiap simpul (\textit{node}) mewakili sebuah entitas tertentu, sedangkan setiap sisi (\textit{edge}) menggambarkan beragam jenis relasi yang mungkin terbentuk di antara entitas-entitas tersebut \cite{Yusrianda2025}. 
Melalui keterhubungan yang bersifat semantik ini, \textit{knowledge graph} memungkinkan penyajian informasi yang lebih kontekstual, terorganisasi, dan dapat diakses secara lebih efisien dibandingkan format teks yang tidak terstruktur.

Sebagai sebuah pendekatan, \textit{knowledge graph} menjadi solusi atas keterbatasan basis data relasional (RDBMS) yang kerap kurang efisien dan fleksibel ketika menangani volume data besar akibat operasi penggabungan (\textit{JOIN}) yang rumit. 
\textit{Knowledge graph} mampu mengakomodasi integrasi data dari berbagai sumber dan format yang kompleks secara lebih efisien, sehingga memudahkan pengguna menganalisis hubungan antarentitas melalui konektivitas informasi yang tervisualisasikan dengan jelas \cite{Prataba_Putra_Widodo_2024}. 
Hal ini sejalan dengan teori \textit{dual-coding}, yang menyatakan bahwa bentuk visualisasi informasi pada \textit{knowledge graph} memungkinkan hubungan antarentitas lebih mudah diproses, diingat, dan dipahami secara lebih mendalam \cite{Prataba_Putra_Widodo_2024}.

\subsection{\textit{Social Network Analysis} (SNA)}
\textit{Social Network Analysis} (SNA) merupakan sebuah pendekatan metodologis yang mempelajari pola hubungan sosial serta interaksi antarindividu, kelompok, organisasi, maupun aktor sosial lainnya, dan secara ringkas dapat dipahami sebagai studi tentang relasi interpersonal dengan memanfaatkan teori graf \cite{GoledzinowskiBlocki2023}. 
Landasan teoretis SNA berakar pada teori graf, di mana sebuah jaringan sosial dimodelkan sebagai graf G = (V, E) dengan V sebagai himpunan simpul (\textit{nodes}) yang melambangkan individu atau entitas, dan E sebagai himpunan sisi (\textit{edges}) yang melambangkan hubungan di antara mereka.
Representasi ini bersifat fleksibel dan dapat berupa graf berarah (\textit{directed}) maupun tidak berarah (\textit{undirected}), berbobot (\textit{weighted}) maupun tidak berbobot, bergantung pada jenis relasi yang dianalisis; misalnya, relasi "\textit{following}" di Twitter bersifat berarah, sementara pertemanan di Facebook bersifat dua arah \cite{Mamatha2024}. 
Dalam konteks penelitian media sosial, SNA dipahami sebagai pendekatan analitis untuk mengidentifikasi struktur sosial sekaligus menjelaskan posisi aktor utama yang berperan penting dalam penyebaran informasi pada suatu jaringan \cite{Sitorus2022}. 
Keunggulan utama metode ini terletak pada kemampuannya mengungkap pola dan dinamika tersembunyi yang tidak tampak pada level individu, seperti bagaimana relasi sosial memengaruhi penyebaran informasi, pembentukan komunitas, serta aliran pengaruh dan kekuasaan dalam jaringan \cite{GoledzinowskiBlocki2023}.

\subsection{Sentralitas}
\label{subsec:sentralitas}
Untuk menentukan aktor yang memiliki peranan penting dalam sebuah jaringan, SNA menggunakan pengukuran sentralitas (\textit{centrality}) yang menjadi inti analisis, karena salah satu kegunaan paling kuat dari teori graf dalam SNA adalah mengidentifikasi individu yang berpengaruh \cite{Mamatha2024}.
Terdapat beberapa ukuran sentralitas yang umum digunakan, yaitu \textit{degree centrality} (yang pada graf berarah terbagi menjadi \textit{in-degree} dan \textit{out-degree centrality}), \textit{eigenvector centrality}, \textit{closeness centrality}, dan \textit{betweenness centrality}. Definisi matematis kelima metrik tersebut mengacu pada kajian teori graf untuk analisis jaringan sosial \cite{Mamatha2024} dan diuraikan sebagai berikut.

\begin{enumerate}
	\item \textbf{\textit{In-Degree Centrality}}

	\textit{In-degree centrality} mengukur jumlah sisi yang masuk ke suatu simpul, dinormalisasi terhadap jumlah maksimum koneksi yang mungkin, sehingga menunjukkan tingkat keterhubungan atau popularitas suatu aktor sebagai tujuan relasi dalam jaringan berarah. Nilai metrik ini dirumuskan pada Persamaan~\eqref{eq:indegree}.
	\begin{equation}
	C_{D}^{in}(v) = \frac{\deg^{in}(v)}{n-1}
	\label{eq:indegree}
	\end{equation}
	Keterangan:
	\begin{itemize}[leftmargin=2em,label=$\cdot$]
		\item $\deg^{in}(v)$ = jumlah sisi masuk ke simpul $v$
		\item $n$ = jumlah total simpul dalam jaringan
	\end{itemize}

	\item \textbf{\textit{Out-Degree Centrality}}

	\textit{Out-degree centrality} mengukur jumlah sisi yang keluar dari suatu simpul, dinormalisasi terhadap jumlah maksimum koneksi yang mungkin, sehingga menunjukkan tingkat keaktifan suatu aktor sebagai sumber relasi dalam jaringan berarah. Nilai metrik ini dirumuskan pada Persamaan~\eqref{eq:outdegree}.
	\begin{equation}
	C_{D}^{out}(v) = \frac{\deg^{out}(v)}{n-1}
	\label{eq:outdegree}
	\end{equation}
	Keterangan:
	\begin{itemize}[leftmargin=2em,label=$\cdot$]
		\item $\deg^{out}(v)$ = jumlah sisi keluar dari simpul $v$
		\item $n$ = jumlah total simpul dalam jaringan
	\end{itemize}

	\item \textbf{\textit{Eigenvector Centrality}}

	\textit{Eigenvector centrality} menilai pentingnya sebuah simpul berdasarkan kualitas koneksinya, yaitu keterhubungan dengan simpul-simpul lain yang juga memiliki pengaruh tinggi, sehingga pendekatan rekursif ini lebih baik dalam memodelkan propagasi pengaruh pada jaringan kompleks \cite{Mamatha2024}. Nilai metrik ini dirumuskan pada Persamaan~\eqref{eq:eigenvector}.
	\begin{equation}
	C_E(v) = \frac{1}{\lambda} \sum_{u \in N(v)} A_{vu} \, C_E(u)
	\label{eq:eigenvector}
	\end{equation}
	Keterangan:
	\begin{itemize}[leftmargin=2em,label=$\cdot$]
		\item $\lambda$ = nilai eigen terbesar dari matriks ketetanggaan
		\item $N(v)$ = himpunan tetangga simpul $v$
		\item $A_{vu}$ = elemen matriks ketetanggaan, bernilai 1 jika terdapat sisi dari $v$ ke $u$
	\end{itemize}

	\item \textbf{\textit{Closeness Centrality}}

	\textit{Closeness centrality} dihitung dari kebalikan rata-rata jarak lintasan terpendek dari suatu simpul ke seluruh simpul lainnya, dan menggambarkan seberapa cepat sebuah aktor dapat menjangkau anggota jaringan yang lain. Nilai metrik ini dirumuskan pada Persamaan~\eqref{eq:closeness}.
	\begin{equation}
	C_C(v) = \frac{n-1}{\displaystyle\sum_{u \neq v} d(v,u)}
	\label{eq:closeness}
	\end{equation}
	Keterangan:
	\begin{itemize}[leftmargin=2em,label=$\cdot$]
		\item $d(v,u)$ = jarak terpendek antara simpul $v$ dan $u$
		\item $n$ = jumlah total simpul dalam jaringan
	\end{itemize}

	\item \textbf{\textit{Betweenness Centrality}}

	\textit{Betweenness centrality} mengidentifikasi seberapa sering suatu simpul berada pada lintasan terpendek (\textit{shortest path}) antara pasangan simpul lain dalam jaringan, sehingga menandai aktor yang berperan sebagai jembatan atau \textit{gatekeeper} dalam aliran informasi. Nilai metrik ini dirumuskan pada Persamaan~\eqref{eq:betweenness}.
	\begin{equation}
	C_B(v) = \sum_{s \neq v \neq t} \frac{\sigma_{st}(v)}{\sigma_{st}}
	\label{eq:betweenness}
	\end{equation}
	Keterangan:
	\begin{itemize}[leftmargin=2em,label=$\cdot$]
		\item $\sigma_{st}$ = jumlah total jalur terpendek dari simpul $s$ ke simpul $t$
		\item $\sigma_{st}(v)$ = jumlah jalur tersebut yang melewati simpul $v$
	\end{itemize}

\end{enumerate}

Melalui kelima pengukuran tersebut, kedudukan dan peran setiap aktor di dalam jaringan dapat diidentifikasi secara sistematis, baik aktor yang menempati posisi sentral, aktor yang berfungsi sebagai penghubung antarkelompok, maupun aktor yang memiliki keterhubungan luas dengan aktor lainnya. Dengan demikian, struktur jaringan beserta pola relasi dan alur transmisi yang terbentuk di dalamnya dapat dianalisis secara lebih mendalam dan menyeluruh.

\subsection{Python}
Python adalah bahasa pemrograman tingkat tinggi yang pertama kali ditemukan oleh Guido van Rossum di Stichting Mathematisch Centrum (CWI), Amsterdam pada tahun 1991, yang pengembangannya terinspirasi oleh bahasa pemrograman ABC \cite{Rahman2023}.
Python merupakan bahasa pemrograman yang menggunakan interpreter untuk menjalankan kode programnya, sehingga kode dapat diterjemahkan secara langsung dan dijalankan di berbagai platform seperti Windows, Linux, macOS, hingga Raspberry Pi. 
Salah satu keunggulan Python dibandingkan bahasa pemrograman lain adalah sifatnya yang open source, sintaksnya yang sederhana dan menyerupai bahasa Inggris sehingga mudah dipahami dan dibaca, serta kemampuannya mengadopsi berbagai paradigma pemrograman, baik prosedural seperti bahasa C, berorientasi objek seperti Java, maupun fungsional seperti Lisp. Karena kemampuannya dalam menangani data besar dan melakukan perhitungan matematika yang kompleks, Python banyak digunakan oleh para programmer dan peneliti untuk berbagai keperluan, seperti pengembangan aplikasi web, aplikasi desktop, \textit{Internet of Things} (IoT), serta proyek \textit{data science} \cite{Rahman2023}.

\subsection{NetworkX}
Dasar teori penggunaan pustaka NetworkX berakar kuat pada implementasi komputasi dari Teori Graf (\textit{Graph Theory}). NetworkX adalah pustaka \textit{open-source} Python yang didedikasikan untuk penciptaan, manipulasi, dan studi struktur jaringan yang kompleks, secara langsung menerjemahkan konsep Analisis Jaringan Sosial (SNA) mengenai simpul (\textit{node}) dan hubungan (\textit{edge}) ke dalam struktur data graf. Kemampuan ini memungkinkan peneliti untuk memodelkan jaringan sosial yang kompleks secara matematis dan algoritmik \cite{PrimFloydSDN2022}.

Keunggulan NetworkX terletak pada efisiensi dan akurasinya; pustaka ini menyediakan serangkaian algoritma sentralitas, termasuk Derajat Sentralitas, yang telah diuji dan dioptimalkan secara ketat. Hal ini memastikan bahwa hasil perhitungan sentralitas yang diperoleh konsisten dan akurat, sejalan dengan definisi formal dalam literatur SNA. Selain itu, NetworkX meningkatkan aksesibilitas data dengan mempermudah proses pemodelan jaringan dari format data mentah yang umum, menjadikannya alat yang sangat ideal dan praktis untuk penelitian kuantitatif berbasis data. Dengan demikian, NetworkX berfungsi sebagai jembatan yang andal antara kerangka teori SNA dan analisis data jaringan yang sebenarnya \cite{PrimFloydSDN2022}.

\subsection{JavaScript}
JavaScript adalah bahasa pemrograman yang berbentuk kumpulan script yang berjalan pada suatu dokumen HTML. JavaScript merupakan bahasa pemrograman web yang bersifat \textit{Client Side Programming Language}, yaitu tipe bahasa pemrograman yang pemrosesannya dilakukan oleh \textit{client}. Aplikasi \textit{client} yang dimaksud merujuk kepada \textit{web browser} seperti Google Chrome, Mozilla Firefox, Opera Mini, dan sebagainya. Dalam pengembangan situs web pemungutan suara, JavaScript digunakan untuk membangun bagian \textit{frontend} situs web, sehingga JavaScript menjadi bahasa dasar dalam proses pengodean (\textit{coding}) sistem yang dirancang \cite{Pramadipta2024}.

\subsection{TypeScript}
TypeScript merupakan bahasa pemrograman berbasis JavaScript yang menambahkan fitur \textit{strong-typing} serta konsep pemrograman berorientasi objek klasik seperti \textit{classes} dan \textit{interfaces}. Dalam dokumentasinya, TypeScript disebut sebagai \textit{super-set} dari JavaScript, yang berarti seluruh kode JavaScript juga merupakan kode TypeScript. Bahasa pemrograman ini menawarkan \textit{classes}, \textit{modules}, dan \textit{interfaces} yang memudahkan pengembang dalam membangun aplikasi yang kompleks, dan hal inilah yang membedakannya dari JavaScript. Adapun keunggulan dan fitur yang dimiliki TypeScript antara lain mendukung \textit{class} dan \textit{module}, \textit{static type-checking}, mendukung fitur ES6, definisi \textit{library} API yang jelas, dukungan bawaan untuk \textit{JavaScript packaging}, kemiripan sintaks dengan \textit{backend}, serta merupakan \textit{superset} dari JavaScript \cite{Permana2022}.

\subsection{React / ReactJS}
React.Js adalah \textit{library front-end} yang dikembangkan oleh Facebook dan sering digunakan sebagai pendukung dari \textit{web-framework}. React.Js memiliki beberapa keunggulan seperti menambah kecepatan (\textit{speed}), kesederhanaan (\textit{simplicity}), dan skalabilitas (\textit{scalability}). React.Js memungkinkan pengembang situs web untuk membangun komponen UI yang lebih interaktif, \textit{stateful}, dan \textit{reusable}. Dalam rancang bangun situs web ini, React.Js digunakan sebagai \textit{framework} dari bahasa pemrograman JavaScript untuk membangun komponen-komponen \textit{frontend} situs web, sehingga proses perancangan dan pembangunan situs web dapat dilakukan dengan lebih mudah \cite{Pramadipta2024}.

\subsection{Teknologi Rendering Graf pada Visualisasi Web}
\label{subsec:teknologi-rendering-graf}
Visualisasi graf berbasis \textit{node-link diagram} di lingkungan web dapat dicapai melalui dua pendekatan yang berbeda secara filosofis: memanfaatkan pustaka visualisasi siap pakai (\textit{high-level libraries}) seperti D3.js, Cytoscape.js, Sigma.js, G6.js, dan ECharts.js yang telah mengenkapsulasi algoritma tata letak dan metode \textit{rendering} ke dalam antarmuka pemrograman aplikasi, atau membangun mesin \textit{rendering} sendiri (\textit{build-your-own}) di atas primitif grafis \textit{native} peramban, yaitu \textit{Scalable Vector Graphics} (SVG) yang bergantung pada \textit{Document Object Model} (DOM) sehingga presisi namun berbiaya \textit{render} tinggi, \textit{Canvas} yang memanipulasi piksel langsung tanpa DOM sehingga lebih ringan untuk animasi, atau \textit{Web Graphics Library} (WebGL) yang memanfaatkan GPU untuk pemrosesan paralel sehingga tercepat pada visualisasi berskala besar \cite{Zhao2025}. Signifikansi perbedaan ini dikonfirmasi secara empiris oleh Zhao \textit{et al.} melalui eksperimen terkontrol pada 481 \textit{dataset} graf berskala 100 hingga 200.000 \textit{node}, yang menunjukkan bahwa kombinasi pustaka berbasis WebGL mampu mempertahankan \textit{frame rate} di atas 30 fps (ambang minimum persepsi visual manusia yang mulus) hingga skala 7.000 \textit{node}, jauh melampaui kombinasi berbasis SVG yang hanya bertahan hingga 600 \textit{node} pada kondisi yang sama \cite{Zhao2025}. Kebutuhan akan visualisasi graf interaktif semacam ini turut ditegaskan pada domain lain yang menghadapi tantangan serupa, yaitu visualisasi ketergantungan kode (\textit{code dependency}) pada rekayasa perangkat lunak; Borowski \textit{et al.} \cite{Borowski2022GraphBuddy} menunjukkan melalui evaluasi pengguna bahwa alat penjelajahan graf interaktif secara signifikan mempermudah pemahaman struktur relasional yang kompleks dibandingkan representasi statis atau tekstual, sekaligus menemukan bahwa penyajian detail yang berlebihan tanpa penyaringan (\textit{filtering}) justru menurunkan kegunaan dan performa visualisasi.

\subsection{Supabase}
Supabase merupakan sebuah platform \textit{Backend-as-a-Service} (BaaS) bersifat \textit{open-source} yang menyediakan serangkaian layanan \textit{backend} secara terintegrasi untuk mempermudah pengembangan aplikasi tanpa perlu membangun infrastruktur server dari awal. Supabase dibangun di atas basis data PostgreSQL yang andal dan menyediakan berbagai fitur utama seperti database relasional, autentikasi pengguna (\textit{authentication}), penyimpanan berkas (\textit{storage}), serta dukungan \textit{real-time} yang memungkinkan sinkronisasi data secara langsung antara server dan aplikasi klien. Dalam perancangan aplikasi pemrosesan input data pada Kumbah Laundry yang dikembangkan menggunakan Flutter, Supabase berperan sebagai \textit{backend} yang menangani proses penyimpanan, pengelolaan, dan pengambilan data pelanggan maupun transaksi laundry. Integrasi antara Flutter sebagai kerangka kerja pengembangan antarmuka (\textit{frontend}) dan Supabase sebagai \textit{backend} dilakukan melalui \textit{REST API} maupun pustaka (\textit{library}) resmi Supabase untuk Flutter, sehingga data yang diinputkan melalui aplikasi dapat langsung tersimpan dan diakses secara aman dan efisien pada basis data. Dengan demikian, penggunaan Supabase memberikan kemudahan dalam pengembangan aplikasi karena menyederhanakan proses pengelolaan basis data, mempercepat waktu pengembangan, serta menjamin keamanan dan ketersediaan data secara \textit{real-time} \cite{Murizar2025}.

\subsection{Pengujian Sistem (\textit{Black-Box Testing})}
\label{subsec:blackbox-testing}
Pengujian merupakan tahap evaluasi yang bertujuan menilai kualitas perangkat lunak berdasarkan indikator yang telah ditentukan, sekaligus berfungsi sebagai tinjauan akhir terhadap pengkodean, desain, dan spesifikasi sistem sebelum digunakan oleh pengguna. \textit{Black-box testing} merupakan salah satu pendekatan pengujian perangkat lunak yang hanya memeriksa hasil keluaran berdasarkan nilai masukan yang diberikan, tanpa memperhatikan struktur maupun kode program internal yang menghasilkan keluaran tersebut \cite{Lubis_Ginting_2024}.

\textit{Black-box testing} bertujuan memverifikasi bahwa setiap fungsi sistem beroperasi sesuai dengan kebutuhan yang telah ditetapkan, sehingga penguji dapat menyusun himpunan kondisi masukan untuk menguji seluruh persyaratan fungsional suatu program, mengidentifikasi kesalahan yang muncul, kemudian melakukan perbaikan hingga sistem dinyatakan layak digunakan. Metode ini juga dikenal sebagai \textit{behavioral testing}, \textit{specification-based testing}, \textit{input/output testing}, atau \textit{functional testing} karena berfokus pada persyaratan fungsional perangkat lunak, bukan pada struktur kode di dalamnya \cite{Lubis_Ginting_2024}.

Kelebihan teknik \textit{black-box testing} terletak pada tidak diperlukannya pengetahuan khusus mengenai bahasa pemrograman tertentu karena pengujian dilaksanakan dari sudut pandang pengguna, sedangkan kelemahannya terletak pada kesulitan menyusun kasus uji tanpa spesifikasi yang jelas serta kemungkinan sebagian bagian \textit{back-end} tidak teruji sama sekali. Metode ini digunakan untuk mendeteksi berbagai kategori kesalahan, antara lain fungsi yang tidak sesuai atau hilang, kesalahan pada struktur data maupun saat mengakses basis data eksternal, kesalahan inisialisasi dan terminasi, serta kesalahan pada antarmuka. Pelaksanaannya diawali dengan mengidentifikasi kebutuhan pengujian, menyusun kasus uji, kemudian menjalankan setiap kasus uji pada aplikasi untuk memverifikasi kesesuaian fungsi, masukan, dan keluaran perangkat lunak dengan spesifikasi yang telah ditetapkan \cite{Lubis_Ginting_2024}.

\subsection{\textit{System Usability Scale} (SUS)}
\label{subsec:teori-sus}
\begin{sloppypar}
Berbeda dengan \textit{black-box testing} yang menilai kebenaran fungsional dari sudut pandang keluaran sistem, \textit{usability} atau kebergunaan menilai sejauh mana suatu sistem dapat digunakan secara efektif, efisien, dan memuaskan oleh pengguna dalam konteks penggunaannya, sehingga pengukurannya bersifat subjektif dan berbasis persepsi pengguna, bukan sekadar benar-salahnya keluaran sistem \cite{Brooke1996}. \textit{System Usability Scale} (SUS) adalah instrumen pengukuran \textit{usability} yang dikembangkan oleh Brooke pada tahun 1986 di Digital Equipment Corporation sebagai metode ``\textit{quick and dirty}'', yaitu murah dan cepat diadministrasikan, namun tetap andal untuk memberikan gambaran umum \textit{usability} suatu sistem \cite{Brooke1996}.
\end{sloppypar}

SUS berbentuk skala Likert 10 item dengan 5 poin (Sangat Tidak Setuju hingga Sangat Setuju), yang item-itemnya dipilih dari kumpulan 50 pernyataan kandidat berdasarkan pernyataan mana yang memicu respons paling ekstrem dan paling konsisten antarresponden pada uji coba awal. Kesepuluh item disusun berselang-seling antara pernyataan bermuatan positif dan negatif, bertujuan mencegah bias respons akibat responden menjawab tanpa benar-benar membaca dan memikirkan tiap pernyataan \cite{Brooke1996}.

\begin{sloppypar}
Penilaian SUS dihitung dengan menjumlahkan kontribusi skor tiap item, yang masing-masing bernilai 0 hingga 4. Untuk item bernomor ganjil (bermuatan positif), kontribusi dihitung sebagai posisi skala dikurangi 1; untuk item bernomor genap (bermuatan negatif), kontribusi dihitung sebagai 5 dikurangi posisi skala. Jumlah seluruh kontribusi kemudian dikalikan 2,5 untuk memperoleh skor akhir pada rentang 0 hingga 100 \cite{Brooke1996}. Skor SUS bukan merupakan persentase, melainkan skor komposit yang hanya bermakna secara keseluruhan; skor tiap item secara individual tidak dapat diinterpretasikan tersendiri \cite{Brooke1996}.
\end{sloppypar}

\section{Analisis Perbandingan Metode}

Berdasarkan Tabel~\ref{tab:ringkasan-pustaka}, studi-studi pada domain dataset dan pemodelan jaringan sanad menunjukkan dua keterbatasan utama. Studi yang berfokus pada pembangunan dataset, seperti Mghari \textit{et al.} (2022) dan Farooqi \textit{et al.} (2024), menghasilkan korpus berskala besar namun tidak menyediakan antarmuka eksplorasi bagi pengguna. Studi yang menerapkan SNA dan knowledge graph, seperti Shuaib \textit{et al.} (2022), Mahmoud \textit{et al.} (2024), dan Hendawi \textit{et al.} (2024), menghasilkan model analitik yang hanya dapat dioperasikan melalui perangkat lunak desktop atau keahlian teknis khusus. Sanadset 650K \cite{MGHARI2022108540} dipilih sebagai data utama penelitian ini karena skalanya yang mencakup 650.986 entri dari 926 kitab melampaui dataset lain, bersifat terbuka, dan strukturnya yang memisahkan komponen matan, sanad, serta perawi memudahkan transformasi ke format knowledge graph.

Berdasarkan Tabel~\ref{tab:visualisasi-perawi}, keempat studi visualisasi perawi bertumpu pada perangkat desktop (Pajek, Gephi) atau sistem web yang terbatas pada rantai perawi tunggal. Ahmad \textit{et al.} \cite{ahmad2025hadith} menerapkan metrik sentralitas dan graf, tetapi hasilnya berupa model analitik, bukan antarmuka yang dapat dijelajahi pengguna. KASHAF \cite{shafie2021kashaf} menyediakan antarmuka web, tetapi fokusnya pada pencarian dan klasifikasi takhreej, bukan eksplorasi posisi sentralitas perawi. Tidak ada studi yang mengintegrasikan knowledge graph perawi dengan metrik sentralitas SNA dalam satu antarmuka web yang dapat diakses tanpa perangkat lunak khusus.

Penelitian ini mengisi celah tersebut dengan membangun \textit{front-end} berbasis web yang menyajikan \textit{knowledge graph} perawi dari \textit{Sanadset 650K} beserta skor sentralitas SNA (\textit{degree} (\textit{in-degree} dan \textit{out-degree}), \textit{betweenness}, \textit{eigenvector centrality}) dalam satu antarmuka yang dapat dijelajahi pengguna tanpa instalasi perangkat lunak khusus.

